 \documentclass[12pt,titlepage]{report}
\usepackage{amsmath}
\usepackage{latexsym}
\usepackage{amsfonts}
\usepackage[normalem]{ulem}
\usepackage{array}
\usepackage{amssymb}
\usepackage{graphicx}
\usepackage[top=1.5cm, bottom=1.5cm, left=1.5cm, right=1.5cm]{geometry}

\usepackage{subfig}
\usepackage{wrapfig}
\usepackage{wasysym}
\usepackage{enumitem}
\usepackage{adjustbox}
\usepackage{ragged2e}
\usepackage[svgnames,table]{xcolor}
\usepackage{tikz}
\usepackage{longtable}
\usepackage{changepage}
\usepackage{setspace}
\usepackage{hhline}
\usepackage{multicol}
\usepackage{tabto}
\usepackage{float}
\usepackage{multirow}
\usepackage{makecell}
\usepackage{fancyhdr}
\usepackage[toc,page]{appendix}
\usepackage[hidelinks]{hyperref}
\usetikzlibrary{shapes.symbols,shapes.geometric,shadows,arrows.meta}
\tikzset{>={Latex[width=1.5mm,length=2mm]}}

\usepackage[utf8]{inputenc}
\usepackage[T1]{fontenc}
\TabPositions{0.5in,1.0in,1.5in,2.0in,2.5in,3.0in,3.5in,4.0in,4.5in,5.0in,5.5in,6.0in,}

\urlstyle{same}



\setcounter{tocdepth}{5}
\setcounter{secnumdepth}{5}



\setlistdepth{9}
\renewlist{enumerate}{enumerate}{9}
		\setlist[enumerate,1]{label=\arabic*)}
		\setlist[enumerate,2]{label=\alph*)}
		\setlist[enumerate,3]{label=(\roman*)}
		\setlist[enumerate,4]{label=(\arabic*)}
		\setlist[enumerate,5]{label=(\Alph*)}
		\setlist[enumerate,6]{label=(\Roman*)}
		\setlist[enumerate,7]{label=\arabic*}
		\setlist[enumerate,8]{label=\alph*}
		\setlist[enumerate,9]{label=\roman*}

\renewlist{itemize}{itemize}{9}
		\setlist[itemize]{label=$\cdot$}
		\setlist[itemize,1]{label=\textbullet}
		\setlist[itemize,2]{label=$\circ$}
		\setlist[itemize,3]{label=$\ast$}
		\setlist[itemize,4]{label=$\dagger$}
		\setlist[itemize,5]{label=$\triangleright$}
		\setlist[itemize,6]{label=$\bigstar$}
		\setlist[itemize,7]{label=$\blacklozenge$}
		\setlist[itemize,8]{label=$\prime$}



\pagestyle{plain}


\renewcommand{\headrulewidth}{0pt}
\setlength{\topsep}{0pt}\setlength{\parindent}{0pt}


\renewcommand{\arraystretch}{1.3}


 \renewcommand{\baselinestretch}{1.5}
\begin{document}
\renewcommand{\thesection}{\arabic{section}} % utilisé pour spécifier la numérotation des sections

\begin{titlepage}
\newcommand{\HRule}{\rule{\linewidth}{0.5mm}}
\center

	\includegraphics[width=0.45\textwidth]{./image2.png}\\[1cm]
   
	\includegraphics[width=0.45\textwidth]{./image1.png}


\HRule \\[0.4cm]
{ \huge \bfseries Cahier des charges \\[0.15cm] }
Classification colorimétrique de nuages de points 3D
\HRule \\[1.5cm]
Ronan Collier,
Mathieu Letrone,
Tri-Thien Truong
\\[1cm]
\today \\ [1cm]
Version 1
\end{titlepage}


\section{Introduction}

Nous tenons avant tout à remercier l'entreprise TagLabs pour avoir participer à la rédaction de ce cahier des charges.

\vspace{\baselineskip}
Taglabs est une jeune entreprise créée il y a deux ans par Yan Koch. L’entreprise s’inscrit dans le domaine de la modélisation 3D d’ouvrages. Ils proposent la modélisation et l’exploitation de nuages de points. Mais, ils travaillent surtout en interne sur un logiciel « ScanSap », le but de ce logiciel est d’exploiter les nuages de points 3D avec efficacité et simplicité inégalées.

Voulant continuer leur développement dans ce domaine encore nouveau, l'entreprise cherche maintenant à améliorer leurs outils, afin de compléter l'exploitation des nuages de points. L'ensemble de ces fonctionnalités permettent à leurs clients de pouvoir analyser un environnement en numérique, à un instant précis (qui sera sous la forme d'un scan de nuages de points). Par exemple, une entreprise peut avoir le besoin d'avoir un scan d'une de leur usine, afin d'analyser le positionnement de leurs machines, les potentielles fuites au niveau des tuyaux, etc.

\vspace{\baselineskip}
L'équipe qui sera à la charge de ce projet est composé de trois étudiants en informatique à Polytech Nantes. Dans le cadre de la quatrième année dans la formation d'ingénieur informatique, nous devons réaliser un projet transversal avec une entreprise.

\vspace{\baselineskip}
\begin{justify}
\textbf{MOA :}
\end{justify}\par


\begin{table}[H]
 			\centering
\begin{tabular}{cccc}
\multicolumn{1}{c}{\Centering \textbf{QUI}} & 
\multicolumn{1}{c}{\Centering \textbf{RÔLE}} & 
\multicolumn{1}{c}{\Centering \textbf{MAIL}} & 
\multicolumn{1}{c}{\Centering \textbf{MOBILE}} \\

\multicolumn{1}{c}{Yan Koch} & 
\multicolumn{1}{c}{Président} & 
\multicolumn{1}{c}{yankoch@taglabs.fr } & 
\multicolumn{1}{c}{0660239733} \\

\multicolumn{1}{c}{Robin Kervadec} & 
\multicolumn{1}{c}{Ingénieur} & 
\multicolumn{1}{c}{rkervadec@taglabs.fr} & 
\multicolumn{1}{c}{0619656021} \\

\end{tabular}
\end{table}

\begin{justify}
\textbf{MOE :}
\end{justify}\par


\begin{table}[H]
 			\centering
\begin{tabular}{cccc}
\multicolumn{1}{c}{\textbf{QUI}} & 
\multicolumn{1}{c}{\textbf{RÔLE}} & 
\multicolumn{1}{c}{\textbf{MAIL}} & 
\multicolumn{1}{c}{\textbf{MOBILE}} \\

\multicolumn{1}{c}{Tri-thien Truong} & 
\multicolumn{1}{c}{Développeur} & 
\multicolumn{1}{c}{tri-thien.truong@etu.univ-nantes.fr} & 
\multicolumn{1}{c}{0631193663} \\

\multicolumn{1}{c}{Ronan Collier} & 
\multicolumn{1}{c}{Développeur} & 
\multicolumn{1}{c}{ronan.collier@etu.univ-nantes.fr} & 
\multicolumn{1}{c}{0666847162\tab } \\

\multicolumn{1}{c}{Mathieu Letrone} & 
\multicolumn{1}{c}{Développeur} & 
\multicolumn{1}{c}{mathieu.letrone@etu.univ-nantes.fr\tab } & 
\multicolumn{1}{c}{0789662916\tab } \\

\end{tabular}
 \end{table}
 
\begin{justify}
\textbf{Tuteur :}
\end{justify}\par 

\begin{table}[H]
 			\centering
\begin{tabular}{cccc}
\multicolumn{1}{c}{\Centering \textbf{QUI}} & 
\multicolumn{1}{c}{\Centering \textbf{RÔLE}} & 
\multicolumn{1}{c}{\Centering \textbf{MAIL}} & 
\multicolumn{1}{c}{\Centering \textbf{MOBILE}} \\

\multicolumn{1}{c}{Nicolas Normand} & 
\multicolumn{1}{c}{Département informatique} & 
\multicolumn{1}{c}{Nicolas.Normand@univ-nantes.fr } & 
\multicolumn{1}{c}{0240683207} \\

\end{tabular}
\end{table}

\section{Modèle du domaine / vocabulaire}

La solution devra permettre d’isoler, classifier un élément dans le nuage de points pour une meilleure visibilité et compréhension. Pour cela, la classification se basera selon la plage de couleur
\begin{itemize}
	\item  s'il est déjà en couleur: on applique le filtre de plage colorimétrique pour isoler les éléments (exemple: isoler les tubes proches du blanc).\par

	\item  s'il est en intensité de gris, on passe d'abord par de la fausse couleur, étape qui permet à l'utilisateur de mieux percevoir le nuage, puis on applique le filtre par plage colorimétrique.\par

\end{itemize}\par

Vocabulaire :



\section{Besoins fonctionnels}

Faire un tableau avec Fonctions, Critères, Niveaux, Flexibilité (aller dans les détails)

Description des acteurs / utilisateurs (situation, expérience, expertise technique, niveau de formation, etc.

Liste hierarchisee des besoins du client (MosCow : Must, Should, Could, Won't)

Utilisez des schémas pour expliquer (fonctionnement, processus, IHM, etc.).

\begin{tabular}{|p{6cm}|p{4cm}|p{3cm}|p{3cm}|}
  \hline
  Fonctions 
  & Critères 
  & Niveaux
  & Flexibilité \\
  \hline
  Isoler un élément dans un nuage de points donné, selon sa plage de couleur 
  & Efficacité de la distinction des couleurs
  & blablabla
  & blablabla\\
  \hline
  Faire apparaître des couleurs sur des nuages de points en intensité de gris
  & blablabla
  & blablabla
  & blablabla\\
  \hline
\end{tabular}

- Lire le scan
- Interpreter les couleurs dans le nuage de points
- Isoler 

\section{Besoins non fonctionnels}


\begin{tabular}{|p{6cm}|p{11cm}|}
    \hline
    Fonctions
    & Commentaires \\
    \hline
    Traitement des problèmes de “mouchetage” 
    & Dans un premier temps, le mouchetage ne sera pas pris en compte dans le traitement du nuage de points.\\
    \hline
    Performances de l'ordinateur
    & A cause des nuages de points, il est nécessaire d'avoir des ordinateurs un minimum performants. Cela veut donc impliquer le fait que le programme développé nécessitera aussi un ordinateur performant. \\
    \hline
    Fichiers en entrée
    & L'utilisation du programme devra prendre en compte la conformité des scans que l'on nous fournira \\
    \hline
    Language utilisé
    & Il n'y a aucun language requis par le client, nous utilisons \\
    \hline
\end{tabular}


\section{Evolution potentielle des besoins}


\section{Limites du projet}

Les fichiers de scans en entrée.
On a des limites sur le traitement du mouchetage (très complexe)

\section{Risques}

Risque par rapport aux distinctions des éléments
Risque par rapport à une mauvaise isolation de l'élément voulu
Risque : mouchetage (visible directement sur le scan)


\section{Plannings}


\section{Gestion de projet}


\section{Connaissances utiles}
\begin{justify}

\end{justify}\par



\end{document}\textbf{\textbf{}}